% !PRINTED 111
\weberpage{111}
%<<CONT

\begin{equation}
 y^2+z^2=\sqrt[n]{b^2+c^2},
\tag{8}
\end{equation}
which determines the absolute value of $x$. Furthermore, from (6) and (7) one obtains
\begin{equation}
 \varphi(y,z)={ (y-iz)^n(b+ic)-(y+iz)^n(b-ic) \over 2i}=0.
\tag{9}
\end{equation}
Now $\varphi(y,z)$ is an integral homogeneous function of the $n$th degree in the two variables $y,z$, and indeed with real coefficients, since replacing $i$ by $-i$ changes nothing in $\varphi(y,z)$. If it is ordered by descending powers of $y$, the first term is $cy^n$, hence, by assumption, different from zero. Putting
\[
 y=\lambda z,
\]
we obtain from (9)
\begin{equation}
 \varphi(\lambda,1)=0,
\tag{10}
\end{equation}
therefore an equation of the $n$th degree for $\lambda$, which by X certainly has a root. Once $\lambda$ is determined, it follows from (8) that
\begin{equation}
 z=\sqrt{\frac{\sqrt[n]{b^2+c^2}}{1+\lambda^2}},\qquad
 y=\lambda\sqrt{\frac{\sqrt[n]{b^2+c^2}}{1+\lambda^2}}.
\tag{11}
\end{equation}
Through this the equations (8), (9) are in fact satisfied; from (9), however, follows
\[
 \frac{(y+iz)^n}{b+ic}=\frac{(y-iz)^n}{b-ic},
\]
and from (8) that the common value of these two fractions is $\pm1$. Therefore we have
\[
 (y+iz)^n=\pm(b+ci),\qquad (y-iz)^n=\pm(b-ci),
\]
and thus the sign of the square root in (11) can be so determined that equations (6) and (7) are fulfilled.

\begin{center}
{\large \S.\ 33.}

\vspace{0.5em}
{\large\gesperrt{Solution of Pure Equations by Trigonometric Functions.}}
\end{center}

The solvability of a pure equation can be shown far more completely and simply if one presupposes the trigonometric functions and their simplest properties as known. These functions are, it is true, foreign to
%CONT>>
